\begin{abstract}
A magic square is a combinatorial paradigm of \emph{balance in inequality}: its $n^2$ entries $\{1,\dots,n^2\}$ are strictly distinct, yet every row, column, and main diagonal sums to the same constant $\mu_n=n(n^2+1)/2$. Lifting the centred grid cells into $\R^3$ as radial vectors $\mathbf p_{ij}(\lambda)=(x_j,y_i,\lambda z_{ij})$ and taking their convex hull—the \emph{Magic Gem} $\mathcal G_\lambda(S)$—raises a central mechanical question: when does combinatorial line-sum equality suffice for physical equilibrium in three dimensions, and when a nonlinear physical realization breaks that balance, how can an opposing inequality restore it? We test this principle across all $880$ order-four magic squares by progressing from one-dimensional radial spokes to three-dimensional rotation, continuous solid and skeletal hulls, and freestanding stacked columns. For radial spokes from the origin $O$, natural vector lengths balance identically, whereas trimming the spokes to equal unit length breaks equilibrium on all $832$ non-associative squares in 3D and on every square in 2D, where the irrational ring-radius ratio $r_2/r_1=\sqrt{5}$ bounds the resultant below by the Pell convergent $9/4\approx\sqrt{5}$ ($18\sqrt{10}/5-8\sqrt{2}\approx0.0705$). By length--weight duality, radial equilibrium is restored either by assigning positive unequal tip weights $\mathbf w^*>\mathbf 0$ on the unit sphere $S^2$ or by scaling spoke lengths as a power law ($r^{\alpha^*}$ in 2D and cubic vertical vanishing $(\mathbf V_3)_z=0$ in 3D). In three-dimensional rotation, equal point masses at the unequal vertices $\mathbf p_{ij}(\lambda)$ make every semi-magic square an axisymmetric top with universal inertia ratio $I_\perp(\lambda)/I_{zz}=\frac12(1+\lambda^2(n^2+1))$ and spherical-top transition $\lambda_{\mathrm{iso}}=(n^2+1)^{-1/2}$, whereas weighting the vertices by their unequal entries $m_{ij}=s_{ij}$ exposes the quadrant quadrupole $Q_{xy}=\sum z_{ij}x_jy_i$, splitting the horizontal principal moments by $2|Q_{xy}|$ and triggering $180^\circ$ Dzhanibekov flips on $824$ squares. Passing from discrete points to the uniform solid hull $\mathcal G_0(S)$ shifts the centre of mass off $O$ on $829$ squares because interior grid points $P_S^{(1)}$ are omitted and hull facets are weighted by unequal simplex volumes; because the outer hull vertices and inner core points pull in strictly opposite directions ($V_0\overline{\mathbf p}_{\mathrm{outer}}+V_1\overline{\mathbf p}_{\mathrm{inner}}=\mathbf 0$), re-engaging the inner core $\mathcal G_1(S)=\conv P_S^{(1)}$ restores balance via two-density core inclusions, two-gauge nested wireframes, or the complete $\binom{n^2}{2}$-rod network $K_{n^2}$ ($\Sigma_{\mathrm{rods}}=\frac{2n^2-3}{3(n^2-1)}\Sigma(\lambda)$), while every boundary hull unfolds into a non-overlapping papercraft net with $f_0-1$ glue tabs. Finally, compressing the vertical aspect to the cubic print scale $\lambda=1/5$ stabilizes $74.5\%$ of all gem facets on a flat table, and alternating any gem $\mathcal G_\lambda(S)$ with its inverted mirror twin $\mathcal G_\lambda(r^2\overline S)=-\mathcal G_\lambda(S)$ cancels opposing facet tilts to lock chiral blocks into a straight vertical tower.
\end{abstract}
